\documentclass[11pt]{article}

\usepackage[margin=0.9in]{geometry}
\usepackage{amsmath,amssymb}
\usepackage{booktabs}
\usepackage{enumitem}
\usepackage{listings}
\usepackage{xcolor}
\usepackage{hyperref}

\definecolor{backcolour}{rgb}{0.95,0.95,0.92}
\definecolor{codeblue}{rgb}{0.05,0.15,0.55}
\definecolor{codegreen}{rgb}{0.0,0.45,0.0}

\lstset{
    backgroundcolor=\color{backcolour},
    basicstyle=\ttfamily\small,
    keywordstyle=\color{codeblue}\bfseries,
    commentstyle=\color{codegreen},
    stringstyle=\color{purple},
    breaklines=true,
    frame=single,
    showstringspaces=false,
    tabsize=4
}

\title{CPSC 250 \\ More on Fibonacci and Complex Numbers}
\author{Edward Brash}
\date{}

\begin{document}

\maketitle

\section{Motivation}

The Fibonacci sequence is surprisingly rich.

Previously, we introduced the sequence and discussed several algorithms for computing it.

In this lecture, we return to Fibonacci from a more mathematical point of view.

This gives us a chance to discuss:

\begin{itemize}
    \item alternate indexing conventions,
    \item negative Fibonacci numbers,
    \item Binet's formula,
    \item non-integer inputs,
    \item complex numbers in Python,
    \item and plotting mathematically interesting functions.
\end{itemize}

\section{The Common Fibonacci Convention}

A common version of the Fibonacci sequence begins:

\[
1, 1, 2, 3, 5, 8, 13, \ldots
\]

with:

\[
f_1 = 1
\]

\[
f_2 = 1
\]

and:

\[
f_n = f_{n-1} + f_{n-2}
\]

for:

\[
n \ge 3
\]

This convention uses natural-number indexing.

\section{The Zero-Indexed Convention}

Earlier, we used a slightly different convention:

\[
0, 1, 1, 2, 3, 5, 8, \ldots
\]

with:

\[
f_0 = 0
\]

\[
f_1 = 1
\]

and:

\[
f_n = f_{n-1} + f_{n-2}
\]

for:

\[
n \ge 2
\]

This convention is common in programming because indexing often begins at zero.

\section{Extending Fibonacci Backward}

A natural question is:

\begin{quote}
Can we extend the Fibonacci sequence to negative indices?
\end{quote}

The answer is yes.

Using the recurrence:

\[
f_n = f_{n-1} + f_{n-2}
\]

we can solve backward.

Since:

\[
f_1 = f_0 + f_{-1}
\]

and:

\[
1 = 0 + f_{-1}
\]

we obtain:

\[
f_{-1} = 1
\]

Next:

\[
f_0 = f_{-1} + f_{-2}
\]

so:

\[
0 = 1 + f_{-2}
\]

and therefore:

\[
f_{-2} = -1
\]

Continuing gives:

\[
\ldots, 89, -55, 34, -21, 13, -8, 5, -3, 2, -1, 1, 0, 1, 1, 2, 3, 5, \ldots
\]

The signs alternate for negative indices.

\section{Ratio Behavior}

For positive indices, the ratio:

\[
\frac{f_{n+1}}{f_n}
\]

approaches:

\[
\phi = \frac{1+\sqrt{5}}{2}
\]

as:

\[
n \rightarrow +\infty
\]

For negative indices, a related limiting behavior occurs:

\[
\frac{f_{n+1}}{f_n}
\rightarrow
-\frac{1}{\phi}
\]

as:

\[
n \rightarrow -\infty
\]

Thus, the golden ratio appears on the positive side, while its negative reciprocal appears on the negative side.

\section{Binet's Formula}

Binet's formula is:

\[
f_n =
\frac{
\phi^n -
\left(-\frac{1}{\phi}\right)^n
}{
\sqrt{5}
}
\]

This formula works beautifully for integer values of \(n\), including negative integers.

\section{A Strange Question}

Once we have a formula, we can ask a strange but interesting question:

\begin{quote}
What happens if \(n\) is not an integer?
\end{quote}

For example:

\[
f_{0.5}
\]

Using Binet's formula:

\[
f_{0.5}
=
\frac{
\phi^{0.5}
-
\left(-\frac{1}{\phi}\right)^{0.5}
}{
\sqrt{5}
}
\]

But:

\[
a^{0.5} = \sqrt{a}
\]

so this requires:

\[
\sqrt{-\frac{1}{\phi}}
\]

which is the square root of a negative number.

This leads naturally to complex numbers.

\section{Complex Numbers}

The imaginary unit is defined by:

\[
i = \sqrt{-1}
\]

A complex number has the form:

\[
z = a + ib
\]

where:

\begin{itemize}
    \item \(a\) is the real part,
    \item \(b\) is the imaginary part.
\end{itemize}

The magnitude of a complex number is:

\[
|z| = \sqrt{a^2 + b^2}
\]

\section{Complex Numbers in Python}

Python uses:

\[
j
\]

instead of \(i\) for the imaginary unit.

This is common in engineering.

Example:

\begin{lstlisting}[language=Python]
import cmath

a = 4
b = 3

z = complex(a, b)

print(f"The real part of z is {z.real}")
print(f"The imaginary part of z is {z.imag}")
\end{lstlisting}

This creates:

\[
z = 4 + 3j
\]

\section{Polar Form}

The module \texttt{cmath} can convert a complex number to polar form.

\begin{lstlisting}[language=Python]
w = cmath.polar(z)

print(f"The size and phase of z are {w}")
\end{lstlisting}

The result contains:

\begin{itemize}
    \item magnitude,
    \item phase angle.
\end{itemize}

The phase angle is given in radians.

\section{Converting Back to Rectangular Form}

We can also convert from polar form back to rectangular form:

\begin{lstlisting}[language=Python]
z3 = cmath.rect(w[0], w[1])

print(z3)
\end{lstlisting}

\section{Plotting Fibonacci for Real Values}

Now we are ready to ask Python to evaluate Binet's formula for many real values of \(n\), not just integers.

The goal is to make a plot of:

\[
f_n
\]

versus:

\[
n
\]

for real values of \(n\) between, for example:

\[
-30 \le n \le 30
\]

\section{Important Practical Note}

For non-integer values of \(n\), Binet's formula may produce complex values.

Therefore, we may need to plot:

\begin{itemize}
    \item the real part,
    \item the imaginary part,
    \item or both.
\end{itemize}

\section{Plotting Plan}

A reasonable plan is:

\begin{enumerate}
    \item Define the Fibonacci function using Binet's formula.
    \item Generate a list or array of \(n\)-values.
    \item Compute the corresponding Fibonacci values.
    \item Plot the normal integer Fibonacci values as distinct points.
    \item Plot the real and imaginary parts for real-valued input.
    \item Draw the horizontal and vertical axes.
\end{enumerate}

\section{Matplotlib Axis Lines}

Matplotlib provides:

\begin{lstlisting}[language=Python]
plt.axvline(x=0)
plt.axhline(y=0)
\end{lstlisting}

These draw:

\begin{itemize}
    \item a vertical line at \(x=0\),
    \item a horizontal line at \(y=0\).
\end{itemize}

\section{Key Takeaways}

\begin{itemize}
    \item Fibonacci numbers can be indexed in more than one way.
    \item The sequence can be extended to negative indices.
    \item Negative-index Fibonacci numbers alternate signs.
    \item Binet's formula works for integer indices.
    \item Extending the formula to non-integer inputs naturally introduces complex numbers.
    \item Python supports complex numbers directly.
    \item The \texttt{cmath} module provides tools for complex arithmetic.
\end{itemize}

\end{document}
